% =====================================================================
%  Paul V. Rubow, GEORG BRANDES AND UNIVERSITY PHILOSOPHY.
%  Danske Studier 1928, pp. 145--162 („Georg Brandes og Universitetsfilosofien“).
%  ENGLISH TRANSLATION, from transcription.tex (Danish), 2026-10-09, made in the main session.
%  Conventions: TRANSLATION-PLAYBOOK.md.  Page numbers of Danske Studier 1928 in the margin; footnotes at the same
%  marks, numbered by page as in the print.  Titles of works stay in Danish.  Quotations Rubow gives in Danish are
%  translated from his Danish; Heine's German stays German.  Terms as in the repo's Nielsen and Brøchner translations:
%  Tro og Viden = faith and knowledge; Uensartethed = heterogeneity; det Antirationelle = the anti-rational.
%  Misprints marked PRINTED AS IS in the Danish are translated as meant.
% =====================================================================
\documentclass[12pt,a4paper,openany]{book}
\usepackage[utf8]{inputenc}
\usepackage[T1]{fontenc}
\usepackage{libertinus}
\usepackage{libertinust1math}
\usepackage{textalpha}
\usepackage{graphicx}
\usepackage[english]{babel}
\usepackage[protrusion=true,expansion=true]{microtype}
\usepackage{setspace}
\usepackage[margin=1.3in]{geometry}
\usepackage{fancyhdr}
\usepackage{hyperref}
\hypersetup{pdftitle={Georg Brandes and University Philosophy}, pdfauthor={Paul V. Rubow}, hidelinks}
\setlength{\headheight}{15pt}
\pagestyle{fancy}
\fancyhf{}
\fancyhead[LE,RO]{\thepage}
\fancyhead[C]{\textit{P. V. Rubow: Georg Brandes and University Philosophy (1928)}}
\setstretch{1.3}
\usepackage{marginnote}
\newcommand{\opage}[1]{\marginnote{\footnotesize\textit{[#1]}}}

\begin{document}

% ===== TEXT (Danske Studier 1928, pp. 145--162) =====
% --- p. 145 ---
\opage{145}\setcounter{footnote}{0}%

\begin{center}GEORG BRANDES AND UNIVERSITY PHILOSOPHY\end{center}

\begin{center}BY\end{center}

\begin{center}PAUL V. RUBOW\end{center}

\begin{center}\rule{3em}{0.4pt}\end{center}

In a small book, \textit{Georg Brandes og hans Lærere} (Copenhagen 1927), I recently tried to give a sketch of his
development, or, I should rather say, of his ``variations'', by letting the great critic pass in review before a
series of more or less significant literary figures who influenced him at various times. Following a method that
common sense and the ever instructive example of Sainte-Beuve both enjoin, I there distinguished between the
influence of the teachers of his youth and the later stimuli. When one deals with an author it is of the utmost
importance to hold him fast during the period in which his personality comes into being, and not for a moment to
let him out of sight, whereas for the later stages, once his mind has become more stationary, one may allow
oneself to fix attention chiefly on certain resting places and turning points.

While in that work, where everything depended on a clear survey, I was fortunate enough, for some sections, to be
able to refer to premises and documentation in an earlier study (see now my \textit{Litterære Studier}, Copenhagen
1928, pp.~38 ff.), there are other parts where I have had to rely to a greater degree on the reader's good faith and
judgment. This is true above all of the first pages, on Georg Brandes's relation to the philosophy of his own
university, which rested on a thorough study of the large literary remains of the two most influential university
professors, Rasmus Nielsen and Hans Brøchner. In what follows I shall try to give a short characterization of
these thinkers and authors --- not, admittedly, an altogether disinterested portrait, for it has throughout been
made with Georg Brandes's writings in mind, and with the greatest attention to the side they turn towards him.

% --- p. 146 ---
\opage{146}\setcounter{footnote}{0}%

\begin{center}I.\end{center}

\emph{Rasmus Nielsen} was, at the time Brandes came to the university, in his fullest strength. He had just then
freed himself from his earlier almost slavish dependence on Søren Kierkegaard's thoughts, and at the same time as
his philosophy of religion began to take on a distinctive form, he sought to reform metaphysics by drawing on the
exact sciences.

The kind of thinking he chiefly cultivated was the concordat. He wanted to unite the most heterogeneous
phenomena; and the formula for the union was that \emph{precisely because}, for example (and first and foremost),
faith and knowledge are absolutely heterogeneous, they can thrive in the same consciousness. If his dialectic was
not profound, it has nonetheless sharpened problems to a high degree. Instead of levelling and mediating he always
sought out new disagreements and exaggerated the oppositions to the point of absurdity. In this he followed in the
footsteps of his oldest master, since Hegel's method is built on the principle of opposition; but he has his
strength precisely where Hegel has his weakness, in that Hegel lays more weight on showing how the one category is
intimated in the other than on showing that the one is self-contradictory unless the other is added.

But he solved the problems only by cutting the knot. Faith and knowledge are absolutely heterogeneous, to begin
with. Next, will and the faculty of knowledge, their organs, are absolutely heterogeneous. Then, finally, a
distinction is drawn between theoretical and practical knowledge, and even here the same heterogeneity is found!
Rasmus Nielsen was under influence from the most various quarters: from Martensen he went over to Søren
Kierkegaard, and Søren Kierkegaard's and Grundtvig's thoughts he later sought to bring under one hat. But beside
this Christian speculation there flourished in him a radical interest in natural science.

A feeling for the special sciences was something surprising in Danish philosophy, so soon after the nature-less
thinking of Martensen and Kierkegaard had entirely occupied people's minds. Rasmus Nielsen held the Hegelian
speculation to have failed, because its concepts were not sufficiently controlled by the individual scientific
disciplines. He wanted a new encyclopaedia of the sciences to replace the Hegelian one, which still seemed to him
burdened with theology and scholasticism. From the end of the fifties he, who until then had thought only about
religious questions, worked
% --- p. 147 ---
\opage{147}\setcounter{footnote}{0}%
as best he could at mathematics, physics, chemistry and biology. After some less successful
attempts\footnote{\textit{Philosophisk Propædeutik i Grundtræk} (1857), \textit{Philosophie og Mathematik} (the same year), \textit{Mathematik og Dialektik} (1859), \textit{Om Begrebet Nøiagtighed} (1860), \textit{Hr. Professor Steens ``Afregning''} (the same year).}
he finally got well under way, and his series \textit{Forelæsninger over philosophisk Propædeutik} (three volumes
from the beginning of the sixties), \textit{Naturphilosophie} (1873) and \textit{Almindelig Videnskabslære} (1880)
show an ever more intimate working of the principles and results of the special sciences into philosophy.

Georg Brandes did not in his young days take much interest in natural inquiry and the exact sciences, and he could
hardly be captivated by Rasmus Nielsen because of his merits in that direction. But it no doubt had an awakening
effect on him to see all the sciences contribute to the building of a philosophical system, as opposed to the
ambiguous ``immanent'' logic of Hegel and Brøchner. And later, when he met positivist philosophy on his way, he was
in certain respects prepared.

By his many-sidedness and his ability to set up oppositions, Rasmus Nielsen came to mean a great deal to many as an
introducer to the problems of philosophy. Georg Brandes himself speaks of him as an awakener, with recognition and
praise. He had, moreover, a spirited written and above all spoken delivery: people would have it that he was the
most eloquent man in Denmark.

He won Brandes by his personality. He possessed a liveliness of mind which seems to have recalled Brandes's own. He
was in constant development; and if he was not very original, he took up in the course of time the thoughts of a
series of significant minds in an independent way. There was in him an eternal appetite for every intellectual
phenomenon, from religious to mathematical problems; the miracles of the New Testament, Darwinism and the defence
question set him vibrating with equal force. For anyone with an eye for the characteristic, he also seems in his
outward appearance to have offered something of interest, as one can read in Brandes's obituary.

In another respect Georg Brandes must have felt drawn to him as to a kindred nature. Rasmus Nielsen was a
considerable literary talent, with even a certain journalistic bent. From his first appearance he had shown an
unerring eye for what was topical and could make an impression. His philosophy
% --- p. 148 ---
\opage{148}\setcounter{footnote}{0}%
came into being by taking up other men's thoughts at the right time. The theory of the relation between faith and
knowledge is ``freely after'' Søren Kierkegaard. If this reworking has little philosophical value, it has on the
other hand some literary value. The slogan of the absolute heterogeneity of faith and knowledge was an excellent
programme for an intellectual movement in the 1860s, a brilliant popularization of Kierkegaard's ideas, which in
their genuine form were not exactly suited to catch on among the bourgeoisie. And it was an able political feat to
bring the followers of Kierkegaard and of Grundtvig closer to one
another.\footnote{Compare Frederik Jungersen: \textit{Dansk Protestantisme ved S. Kjerkegård, N. F. S. Grundtvig og R. Nielsen}, Copenhagen 1873.}
It was also useful to work at removing religious suspicion from scientific research. By his struggle of many years
to tear theology down and, on the other hand, to preserve Protestant Christianity and the Church, he is remarkable,
rather as a leader of the people than as a critical thinker.

If one wants to see to what degree he was an advocate in his thinking, one should read some of his rather numerous
polemical writings. Here his standpoint in the philosophy of religion changes constantly with his development and
with the various opponents he faces. Nowhere does the tough peasant nature show itself better than in his stubborn
and not wholly unfruitful philosophical polemic. What tone is to be struck, how expansively one is to proceed, all
depends on the opponent he has to deal with. Bishop Martensen he is always ready to give full measure, up to a
hundred and fifty pages; Provost Zeuthen gets less at once; Brøchner, his colleague in the very narrowest sense,
whose very bitter attack cannot be left unanswered, is dispatched in a good three sheets; Georg Brandes and A.~C.
Larsen, who are only young graduates, receive no answer at all.

He is no critical talent; the detailed logical discussion of the argument he avoids, and keeps to certain points
in the opponent that can bring public opinion over to his side: in Martensen the theological universal monarchy is
made ridiculous, in Brøchner the incomprehensible terminology; on Martensen the suspicion is cast that he wants to
put an end to free inquiry, Brøchner is accused of personal enmity and hatred of Christianity. Nothing prevents him
from going
% --- p. 149 ---
\opage{149}\setcounter{footnote}{0}%
round the main issue; for example, ``the absolute heterogeneity'' is passed over in the piece against Brøchner.

Rasmus Nielsen was, finally, something of a stylist. His manner of writing is just as little original as his
thinking. It descends directly from that of Kierkegaard's philosophical writings: the same overloading with climaxes,
repetitions, parallels, antitheses, puns, the whole barbaric ornament. But he gradually succeeded in transforming
the master's passionate and individual delivery into a broad, vivid art of persuasion. Often he abuses the figure of
antithesis to the point of disgust; and he plays a game with the anaphoric sentence that it takes strong nerves not
to be irritated by. In \textit{Om den gode Villie} he lets each of five sections, over two pages, begin with the
question: ``But how is the dualism to be overcome?''\footnote{In part slightly varied; see the work cited, pp.~73 ff.}.
It cannot be denied that his manner of expression gains in pedagogy what it loses in taste. Whatever he touches
comes alive. He knows how to present the most abstruse \textit{metaphysica} in a stimulating way and with a kind of
intelligibility, in so far as what is obscurely thought can ever become what is clearly said. In this he is quite
different from Brøchner, who smothers all life in his logical-terminological armour, and whose intricate sentence
structure seems to have no other principle than to separate what naturally belongs together for the understanding.

Rasmus Nielsen shuns no means of holding the attention captive: popular sayings, bold allusions, familiar
quotations, coarse images in abundance. He even undertakes to renew the technical terms; and when he uses the
forbidding foreign words, he often personifies them a little, which is undeniably more helpful to intuitive
understanding than to the logical process. On the whole he likes to make concepts into persons or tangible things:
``The theology of Scripture, so youthfully fresh, so trusting, so proud of the Bible in its Paradise --- how has it,
after the Fall through the logical-physical, gradually been weakened, bowed, bleached, humbled!''\footnote{The work cited, pp.~48 f.}.
He is fond of exclamations, dialogisms and outright dialogues, not only between A and B and Caius, but also, for
example, between Theology and the Congregation:

\medskip

\begin{quote}\small

``What have you done'', the Congregation asks Theology, sinful by its Fall, ``with my Christmas Gospel?'' Alas,
dear Congregation, do not be
% --- p. 150 ---
\opage{150}\setcounter{footnote}{0}%
angry; criticism has unfortunately torn the Christmas Gospel clean to pieces. But do not lose heart! remember that
it is knowledge that is to help faith; remember what Bishop Martensen says with Anselm: ``I do not venture, O Lord,
to fathom thy depths;'' but I believe that knowledge is to help faith. Do not weep; the Gospel that criticism has
torn to pieces, apologetics can, after all, put together again \dots{}

You look so darkly at me. What can I do about science being objective? Who would have thought, too, that the
historical facts in the Christmas Gospel were so strangely brittle, the natural elements just as brittle as the
supernatural \dots{}\footnote{Ib. p.~49.}

\medskip

\end{quote}

When his wit does not run away with him too much, he also knows how to build up a book simply and in a way all can
understand; this is especially the case with his popular writings after 1870, but also with earlier works. It holds,
for example, already of \textit{En undersøgende Anmeldelse} (1849) and of his most forceful polemic, the
previously cited \textit{Om den gode Villie} (1867), which simply and clearly establishes the impossibility of
theology's supremacy over the other sciences by showing first history's, then the natural sciences', and finally
philosophy's untenable position in such a union. The double work \textit{Religionsphilosophie} (1869) ---
\textit{Naturphilosophie} (1873) is also cleverly planned, \textit{Grundideernes Logik} (1864--66), on the other
hand, confused.

In keeping with this, Georg Brandes is seen in his small reviews from the end of the sixties to take pleasure in
Rasmus Nielsen as an author. In the review of \textit{Om den gode Villie} he permits himself small remarks on the
literary character of the work: ``a broad and racy, a mocking and biting little book, written with an overflowing
delight in production and full of good ideas. Its merriment has nonetheless in places a more Dutch than Danish
character''\footnote{Georg Brandes, \textit{Samlede Skrifter} XIII p.~85.}, ``the excellently witty little piece
that ends with the words \dots{}''\footnote{Ib. 87.}. Brandes's whole account of Rasmus Nielsen's relation to
Martensen and theology belongs in the same place.

When Brandes's former opponent, Rasmus Nielsen's apostate disciple P.~S.~V. Heegaard, wrote a treatise against the
doctrine of his once idolized master, Brandes had the triumph of being able to take Heegaard to task for his lack
of recognition of ``his so highly gifted teacher'' and to weave in some laudatory characterizations: ``he has the
luxuriantly fertile mind, the inspired eye, the rare originality, which are just as essential faculties in a
% --- p. 151 ---
\opage{151}\setcounter{footnote}{0}%
thinker as a strictly scientific training of thought and a keen, reliable understanding''\footnote{Ib. p.~108.}.

It is hardly unwarranted to find some dependence on Rasmus Nielsen's style in these reviews and in the pamphlet
\textit{Dualismen i vor nyeste Philosophi}. % Danish: „Afhængighed at“, PRINTED AS IS for af
One does not arrive, by the mere study of Taine and with Gabriel Sibbern's assistance\footnote{See my
\textit{Litterære Studier} (Copenhagen 1928), pp.~45 f.}, at writing a modern prose that all can understand. At
the least, Georg Brandes's inclination towards dialogue in these works smacks just as much of Rasmus Nielsen as of
Taine: ``[it was] good policy on the part of the philosophical champion to proceed by attack and answer Bishop
Martensen's polite `Good day' with a terrible `Axe-handle'\,''\footnote{G.~B. \textit{Saml. Skr.} XIII p.~91.}.
``Look at \textit{Grundideernes Logik}, says the author [i.e.\ Heegaard]: the best thing would be for Professor
Nielsen, who is after all unable to finish building it, to let it stand at once as an `interesting ruin', that is,
as a kind of Marble Church. `Marble Church yourself', answers Professor Nielsen, calmly pointing to the first part
of the far-famed \textit{Indledning til den rationelle Etik}''\footnote{Ib. p.~107.} (in this work Heegaard, as is
well known, was still at the Nielsenian stage, and it was notoriously never continued). Here there are jokes in an
unmistakably Nielsenian manner: ``Professor Nielsen has indeed one man on guard, but he weeps --- with one eye and
laughs with the other; his right hand knows not what the left is doing. He moves forward on two heterogeneous legs,
a leg of flesh and a wooden leg, on two absolutely split legs, cleft like one of Uffe's Saxons. His head is a
Janus head: with one mouth he proclaims the higher fulfilment of the law of nature through the miracle; with his
other throat he gives cultural consciousness the honour and sings a well-meant hymn of praise to David
Strauss''\footnote{Ib. p.~91 --- as to the substance of the matter there is nothing exaggerated: those who today
hold that Jesus never existed as a historical and actual person can without risk put Rasmus Nielsen on the list of
those who think as they do.}.

How deep the impression of Rasmus Nielsen goes is seen clearly, finally, from the obituary Brandes wrote of
him\footnote{\textit{Saml. Skr.} II pp.~419 ff.} (1884). Here he gives him quite certainly all he deserved as a
personality, as a thinker a little more, and brings out finely and surely his overlooked merits as a writer.

As we have seen, then, he stood in some debt of gratitude to his old professor of the \textit{filosofikum}. He was
reckoned by contemporaries from
% --- p. 152 ---
\opage{152}\setcounter{footnote}{0}%
every camp, with good reason, as one of Rasmus Nielsen's most dangerous enemies. This earned him civilities from
Nielsen's older opponents and drew curses upon him from his followers. Thus he is mentioned with honour in
Martensen's \textit{Om Tro og Viden} (1867) --- whose diction, like that of almost every contribution to this
famous controversy, is a little infected by Rasmus Nielsen's --- as the ``gifted author'' of whom one ``can only
wish that he may not too early get beyond the Platonic \textit{θαυμαζεῖν} (wonder) as the beginning of
philosophy''\footnote{The work cited, p.~8.}. Conversely, Brandes gets a lecture and a scolding in the
aforementioned book by Fr. Jungersen, \textit{Dansk Protestantisme}, where the pamphlet on \textit{Dualismen} is
called ``an extraordinarily frivolous criticism''\footnote{The work cited, pp.~210 f.}. The many times Brandes
himself spoke of Rasmus Nielsen, in writings or in interviews, it was always with high regard, at times almost with
tenderness. If one studies the relation at close range, one sees that more than one thread leads from the great
Voltairean critic back to the Kierkegaardian sophist.

\begin{center}II.\end{center}

The relation to \emph{Hans Brøchner} is far simpler and more transparent. The man to whom Georg Brandes dedicated
his principal work was his real, strict educator; he plays a role in his life similar to that of Vacherot in
Taine's, % Danish: „har spiller“, PRINTED AS IS for han
and, like Vacherot, was no uncritical spectator of his pupil's rapid
development.\footnote{In \textit{Tilskueren} for February 1912, p.~213, Dr.~Poul Levin published Brøchner's
assessment of the dissertation Brandes submitted for the doctorate, which I here take the liberty of reprinting:
``It seems to me that Cand. Brandes's dissertation, judged from a purely philosophical point of view, has two not
inconsiderable shortcomings. The treatment of the problems it raises is somewhat aphoristic, and in showing Taine's
relation to his historical presuppositions in the field of philosophy there is an incompleteness, which means that
it does not become sufficiently evident to what striking degree he lacks originality, depth and consistency as a
philosopher. But these shortcomings of the dissertation are no doubt amply outweighed by its great merits: by its
rich content, which testifies to a comprehensive and well-assimilated wealth of knowledge in the author, by the
keen, fine and spiritedly individualized understanding of its subject, and by the life and freshness of the
presentation. I therefore have no doubt that the dissertation ought to be accepted.''}
In his \textit{Levned} Brandes has piously recorded the day of his first visit to the beloved teacher: ``it was the
2nd of September 1861'', and has given an account in every detail of their intellectual companionship. Here I add
only what a commentator might further have to recall.

Brøchner was a quiet, introverted thinker and scholar, who
% --- p. 153 ---
\opage{153}\setcounter{footnote}{0}%
at that time had not yet published his most important works. But he had them finished in his head and could later
publish them in definitive form in a surprisingly short time. Through lectures and private conversation their
content was known to Brandes long before they appeared in print.

In a private letter to Chr.~K.~F. Molbech Brøchner expressed the wish to leave behind two works: one that could show
he had known something, one that could testify that he had understood how to think.

The first he has displayed in his works on the history of philosophy, where on the one hand --- with the assurance
of an old general-staff officer --- he skeletonizes the most representative metaphysical edifices from Greek
philosophy in its tragic age down to Hegel, and on the other hand reads speculatively the work of the world spirit
in the laws which he believes determine the development of thought, and which thinkers unconsciously follow. A
literary historian naturally has much to object to in these works of learning because of their stiff-necked
deductive procedure --- Brøchner had a predilection for writings in the form of handbooks and all-embracing systems,
and an unparalleled contempt for the connection between philosophy and the other sciences, which makes parts of his
writings, for example the treatment of French materialism in the eighteenth century, altogether misleading; and his
inclination to establish speculative connections weakens the significance of the sections where he was best
grounded, philologically too, as in the account of the philosophy of the
Greeks.\footnote{His insight into medieval thought has benefited others only through the commentary to
Chr.~K.~F. Molbech's translation of Dante (1851--62), for which he supplied the philosophical material.}
But his historical works, notably the handbook in two volumes (1873--74), still command respect by their precision
and thoroughness.

His capacity for independent thought he has proved in the two books that give his own philosophy:
\textit{Problemet om Tro og Viden} (1868) and \textit{Om det Religiøse i dets Enhed med det Humane} (1869).

He begins --- not without a sense of the topical --- with a criticism of contemporary Danish philosophy's and
theology's attempts to solve the question of faith versus knowledge. He rejects the theological philosophy of
faith, represented by Martensen, which wants to build up a supra-logical science of things that it itself admits
cannot be known. He rejects further, with reverence and admiration, Søren Kierkegaard, who breaks the connection
between the personality and
% --- p. 154 ---
\opage{154}\setcounter{footnote}{0}%
universally valid knowledge by defining faith as ``the objective uncertainty held fast in the appropriation of the
most passionate inwardness'', and who, by denying the natural side of man, sets an insoluble disharmony in the
personality and in existence. He rejects, finally, Rasmus Nielsen's doctrine of absolute heterogeneity, as a sham
solution, an anarchy of irreconcilable contradictions.

If one wants to explain the relation between faith and knowledge, one must, according to Brøchner, take one's
point of departure in one of these spheres. And he chooses --- like Hegel --- knowledge, because only knowledge can
give an independent view of existence. Faith could not stand alone, but would have to lean on knowledge in the
face of such a task.

He now builds up a whole monistic and idealistic metaphysics. He rejects once and for all the atomistic way of
looking at things and will operate only with totalities. The whole of existence is an absolute totality, and within
it the individual is in turn a relative totality, which has its task to fulfil within the absolute.

The principle of existence is ideal. The real is a transitional stage in the development of the Idea; the Idea is
both the ground of the totality and its end. Brøchner dissolves all dualisms: knowledge and will are not
heterogeneous, the will has a knowledge as its content, and it develops ``from the first immediate, single
determination by impulse'' through reflection ``to a determination by consciously formed principles, maxims, from
which the unity of a view of life can then
emerge''\footnote{\textit{Om det Religiøse} pp.~18 f.}. Nor is there any difference of quality between
theoretical and practical knowledge. The difference is placed in this, that practical knowledge is personal,
theoretical knowledge impersonal --- but ``the act of knowing, which is subject to the norm of the idea of truth,
has in that norm an `ought', in relation to which there is consciousness of
responsibility''\footnote{Ib. p.~15.}.

Brøchner's philosophy of religion is determined by these basic concepts. The religious relation is seen as the
human spirit's ``absolute relation to the Absolute'', i.e.\ its undivided and unconditional devotion to and
subordination under the eternal and unchanging ground-force of existence. Both knowledge and will must take part.
Knowledge is continued in a postulate: faith. There is in all lower and higher knowledge a moment of faith, namely
trust in the unity and lawfulness of existence, and a moment of devotion: ``In all true knowledge there is a
purification of thought, a giving up of individual prejudice, opinion and bias, in a word: of the selfishness
% --- p. 155 ---
\opage{155}\setcounter{footnote}{0}%
of thought; there is in it a self-surrender to the absolute power of truth''\footnote{Ib. pp.~50 f.}.
The faith that is present in knowledge is rational, and must therefore reject everything that conflicts with the
principle of sufficient reason and of necessity, for example the miracle. The will is continued in self-surrender,
as the individual becomes conscious of itself ``as an organic member of the totality of existence'' and ``chooses to
act accordingly''. The moment of faith shows itself in the ethical, which is ``the expression of the true,
essential willing and acting''; it comes forward ``in every ethical decision in which the individual resolves on
action despite the incompleteness of knowledge and the uncertainty of its realization, of the
outcome''\footnote{Ib. p.~55.}. The moment of devotion is seen in this, that the principle of the ethical is love.

The strife that constantly rages between representatives of knowledge and of faith in its positive form is
therefore not due to an opposition of principle: the cause must lie in the finitude of the human spirit's
conditions. If we disregard everything finite in the positive religions, we find that the most important, or the one
thing needful, in them is the demand for subordination, self-surrender, which is thus also found in knowledge. Faith
and knowledge, then, do not conflict: there remains ``no doubt something like a gap between faith and knowledge,
but a gap that must be able to be traversed by knowledge''\footnote{Ib. p.~51.}. Brøchner believed he could observe
in the history of religions a constant progress towards a more purified religiosity. The highest standpoint is the
religion of humanity, which does not need to be created but only to be pointed out, for the religious is already
present, ``if often unclearly and half-consciously'', in the humane consciousness of our time.

In connection with his religious view Brøchner treats the ethical problems in a very speculative spirit. Free will
is seen in man's capacity to assert himself, according to his infinite determination of essence, over against the
external and accidental. Evil is a disturbance in ``the normal relation between man's being as a relative totality
and as a member of the universal totality, and between the natural and the spiritual side in it.'' The will is evil
when it is self-will and sets itself up against the demand of the Absolute. Striking is the optimism of this
ethical view: wide possibilities are left open for repentance and reconciliation, and distrust is expressed of the
evil will's capacity to endure. Evil is rather an accident, which moreover has the greatest probability of
correcting itself. The will to evil has a natural tendency to turn round into the good;
% --- p. 156 ---
\opage{156}\setcounter{footnote}{0}%
it borrows its strength from the Absolute\footnote{Compare Georg Brandes's \textit{Æsthetiske Studier} p.~116:
``the power of the bad is the strength it has sucked from the good.''}, and so must fight against the ground of its
own being and in the end destroy itself.

When all is said and done, his view is a religious optimism. Values are preserved; nothing essential is lost. He
denies, for example, personal immortality, but holds that the individual spirit, at the destruction of the organism,
is taken up ``as a moment in the life of the Eternal''; it is transformed into ``a moment in God's knowledge of
himself in his eternal kingdom of spirit''\footnote{\textit{Om det Religiøse} p.~189.}. Existence is a cosmos.

Brøchner reproached the Hegelian philosophy for not letting reality come sufficiently into its own. The same can,
however, be said of his own system. It is spiritualistic and transcendent like Hegel's. It is without relation to
the special sciences, and in that respect far less modern than Rasmus Nielsen's. With all its consistency and
firmness it is extremely exclusive, completely closed on all sides. It is hard to know what advantage philosophy
could draw in the future from natural inquiry and exact science, although Brøchner himself pointed the philosophers
of the future to this field of work in his writings on the history of philosophy. The method, too, is the old
``immanent'' development of concepts, however much Brøchner himself inveighed against
it\footnote{See his \textit{Bidrag til Opfattelsen af Philosophiens historiske Udvikling} (1869) pp.~265 f.}.

The above meagre account of Brøchner's religious free thought --- a last shoot on the stem of romantic philosophy,
the speculative testament of an epigone firm in his principles --- may serve as a background to, and to complete,
the interesting account in Brandes's \textit{Levned} of how he became a freethinker, and may be compared with the
summary presentation in my aforementioned study \textit{Georg Brandes og hans Lærere}, pp.~6 ff. It goes without
saying that the intimate influence from master to pupil in the time before Brøchner began his real, central
authorship, and before Brandes published anything at all, is a thing that cannot be made the object of objective
investigation and matter-of-fact discussion.

But one sees at least from Brandes's own words and from Brøchner's printed writings what sort of education the
future critic of naturalism received. His youthful thinking was schooled in an abstract and idealistic philosophy
of concepts. And the thinkers to whom Brøchner introduced him were of the same
% --- p. 157 ---
\opage{157}\setcounter{footnote}{0}%
kind. But regardless of everything that Brandes later let drop as at bottom incompatible with his nature, there
are features in Brøchner's philosophy that were bound to win his full assent. It is, as has been shown, strictly
monistic, and by its unconditional holding fast to one principle it counteracted all the temptations to a
dualistic view of life, of which the age had plenty. Next, it is markedly intellectualist. It asserts the primacy
of knowledge in the human spirit, and chooses knowledge to give the key to the riddles of existence. Finally, it is
avowedly freethinking. Brøchner was a stubborn champion of the unlimited rights of ``free thought'' and was hated as
an enemy of Christianity, an accusation that he rejects with great force and dignity in his \textit{Svar til
Professor R. Nielsen} (1868). His opinion of Christianity can be seen from his view of positive religions in
general; moreover, his book \textit{Om det Religiøse} contains here and there sharp criticism of the chief Christian
conceptions: the miracle, creation, immortality, and of the tendency of Christian doctrine to deny the natural ---
he held to the correctness of Kierkegaard's picture of Christianity against Martensen's. As a young man he had had
to suffer a refusal from the theological degree examination for his openly avowed view of the miracles as mythical.
He was one of those who had introduced Strauss\footnote{\textit{Den christelige Troeslære \dots{} Oversat af
H. Brøchner, Cand. phil.} 2 vols. 1842--43.} and studied Feuerbach, a thinker who has meant much for the
development of the religious standpoint to which he held unshakably.

The writings of Georg Brandes in which one would most naturally look for special and tangible influences from
Brøchner are \textit{Den franske Æstetik i vore Dage}, his most academic work, and \textit{Emigrantlitteraturen},
his freethinker's manifesto. But what is to be found here is microscopic.

The dissertation, in which he so to speak lays out his younger teacher's doctrines before his older teacher, is
certainly a clean piece of pupil's work, which follows the master's method of preparing a philosophical system:
give its presuppositions and pass judgment (on which, as will be seen from the note above, p.~152, the two parties
were nevertheless not agreed). But in the details one finds only odds and ends. The most important passage is
p.~241, where I find influence from Brøchner's \textit{Om det Religiøse}: ``Taine has started from Spinoza's
standpoint, that man is a part of the whole and not `a state within the state'; but
% --- p. 158 ---
\opage{158}\setcounter{footnote}{0}%
he has gone so far as to conclude that man cannot be a citizen of this state either, is only a product without
personality, without freedom and without real rights. He thus comes to ignore self-consciousness for the good
reason that he cannot explain it, since he lacks any concept of \textit{man as a relative totality} and understands
nature from the idea of mechanism instead of from that of organism.''

In the first three volumes of \textit{Hovedstrømningerne}, too, Brandes wanted to continue Brøchner's work; this
was clear to contemporaries, and his opponents hoped to hit the master when they shot at the disciple.
\textit{Emigrantlitteraturen}\footnote{On whose academic character and whole position I may refer to my
\textit{Litterære Studier}, pp.~80 ff.} is dedicated to Brøchner ``with gratitude and friendship''. \textit{Den
romantiske Skole i Tyskland} has the following afterword, not reprinted later:

\medskip

\begin{quote}\small

Dear Professor,

\medskip

It goes without saying that this volume, just as much as the preceding one, is dedicated to you. With my
university lectures as its basis it contains a good deal that could not find room, or would not have been in its
place, in the spoken lecture. It has been my hope while working it out that men such as you would find pleasure in
this book. For me a book is an act. You will judge it as such.

You know Heine's words, Professor: ``Sonderbar! und immer ist es die Religion, und immer die Moral, und immer der
Patriotismus, womit alle schlechten Subjecte ihre Angriffe beschönigen! sie greifen uns an, nicht aus schäbigen
Privatinteressen, nicht aus Schriftstellerneid, nicht aus angeborenem Knechtsinn, sondern um den lieben Gott, um
die guten Sitten und das Vaterland zu retten.'' --- Happily, for me a whole bundle of anonymous articles and
pamphlets is outweighed by one word from you.
\par\hfill Copenhagen, April 1873.\par\hfill Yours\par\hfill G.~B.

\medskip

\end{quote}

But although \textit{Hovedstrømningerne} in its original design is a continuation of Brøchner's polemic with
theologians and dualists --- the very title and the often recurring slogan ``free thought'' come from him --- I do
not believe it will be possible to demonstrate Brøchnerian influences in it in detail. It is only in the sixth
volume, the book on \textit{det unge Tyskland}, that the already ageing Brandes ties the thread to his Hegelian
youth and describes the Young Hegelian struggle for freedom of thought and against
% --- p. 159 ---
\opage{159}\setcounter{footnote}{0}%
clericalism. Here he puts in his teacher's Left-Hegelian interpretation of Hegel's
metaphysics\footnote{Compare again \textit{Georg Brandes og hans Lærere}, pp.~13 ff.}.

There is only a quite small group of Brandes's writings in which one can say that the influence from Brøchner goes
more into particulars: the polemical works from the years on either side of 1870. Moreover, it is possible that the
influence was mutual; at any rate Brøchner's oldest occasional piece, a contribution to the baptismal controversy at
the beginning of the forties, stands far behind those with which he contributed to the dispute over faith and
knowledge. It was Georg Brandes himself who forced the pen into his hand\footnote{See Brandes's \textit{Levned} I
p.~177.} and got him first to criticize his contemporaries' views, and then to develop his own; after that the
rest followed of itself.

In \textit{Dualismen i vor nyeste Philosophi} (1866), which is older than Brøchner's contributions to the
controversy, one should accordingly not expect formal influence, but the line of thought is quite certainly
inspired through and through by Brøchner's; it also agrees excellently with Brøchner's own deductions two years
later in \textit{Problemet om Tro og Viden}. First there is an account of everything tempting in Rasmus Nielsen's
philosophy: suspicion is even cast on the possible opponent: ``the individual who appears not for the sake of the
cause, but in order to bring his own little self into the foreground by means of the great
cause''\footnote{The work cited, p.~11.}. But then ``the medal is turned'', and the explosives are thrown down
with all the greater force. And after confidence in the system has been shaken, it is proved with strict logic both
that the premises of the thesis put forward do not agree with the conclusion, and that the conclusion is in itself
a self-contradiction. Next it is shown that the auxiliary concept ``the anti-rational'' is a hopelessly
obscurantist notion. Finally, the ruinous consequences of the asserted dualism throughout its originator's system
are demonstrated, and the failed attempts at soldering are laid bare, until the system is at last condemned in its
pitifulness as a piece of failed eclecticism.

The basic thoughts are the following. 1\textsuperscript{o} --- Propositions of faith and propositions of
knowledge, as Rasmus Nielsen sets them up, are not absolutely heterogeneous. Propositions such as ``The world was
created in six days'' --- ``The world was not created in six days'' are merely conflicting, incompatible. But even if
faith and knowledge really were heterogeneous, it would still not help, since heterogeneity is only another
expression for incompatibility. And it
% --- p. 160 ---
\opage{160}\setcounter{footnote}{0}%
would presuppose a split both in consciousness and in existence. 2\textsuperscript{o} --- Postulates of faith can
be criticized by knowledge: ``whoever --- be it as believer or as knower --- assumes anything whatever about the
origin of the earth has an assumption that geology can correct''\footnote{Ib. p.~31.}. 3\textsuperscript{o} --- The
``anti-rational'', Rasmus Nielsen's unparalleled discovery, falls entirely outside scientific discussion, since it
cannot be expressed in rational concepts or even in language.

Religion is not criticized in the little work. One remark leaves open the possibility of a religion without
miracles: ``If `religion presupposes a revelation and revelation a miracle, but science denies the miracle and
rejects revelation' \dots{} then the form of the religious that can be reconciled with science can only be one that
is without miracles and without revelation''\footnote{Ib. p.~60. Cf. the passage on Parker, ib. pp.~73 f.}.

\textit{Forklaring og Forsvar, En Antikritik} (1872) is, conversely, influenced in form, but hardly in content, by
Brøchner's intervening polemical writings, notably by \textit{Et Svar til Professor R. Nielsen}. Brandes admired in
the highest degree Brøchner's ability as a polemicist. In his essay on the beloved teacher from 1886 he writes of
it:

\medskip

\begin{quote}\small

His war of attack on the theologians, or on Nielsen's masked-theological philosophy, was a model of superior
tactics. He would usually first develop with the greatest precision the view he meant to dispute, then set out
point by point the proofs of its untenability, then adduce everything that could conceivably be urged from the
opponent's side against these proofs, and explain why this supposed defence would mean nothing; then he proved that
no defence other than the one adduced could be made at all.

When this was done, he dissolved the opponent's conceptual structure, made him concessions, supposed \textit{per
impossibile} that one of the overthrown main points was a firm fortress, and showed that even so the most decisive
determinations of concept contained contradiction upon contradiction. Step by step he went round his opponent's
doctrinal edifice, and, well acquainted with his manner of fighting, he stopped every hole through which he might
try to escape, until he had him safely caught. Then he would usually end by enumerating the main points in all the
contradictions shown, and closed at last the criticism, held in the most serious tone, with an irony that one could
in all truth, and without abuse of the word, call terrible.\footnote{\textit{Saml. Skr.} II pp.~429 f.}

\medskip

\end{quote}

% --- p. 161 ---
\opage{161}\setcounter{footnote}{0}%
Brandes's own tactics in the five short chapters of his \textit{Forklaring} of 1872 are not much different from
those he has described here. It is especially clear that he imitates Brøchner every time he passes from the
defensive to the offensive. He first meets the numerous attacks in the press. By deftly chosen quotations the
opponents are made to contradict themselves every time they put forward a claim; and he is not afraid to use his
victory in such a way that he grants the opponent to be right in what has just been proved to be wrong, in order
then to show that it does not help him anyway:

\medskip

\begin{quote}\small

I have said: That a literature lives in our day is shown by its putting problems up for debate. I have said, which
\textit{Dagbladet} passes over without further ado: in our day; I have not claimed that in the time of Homer or
Hesiod the proof of a literature's life lay in its discussion of problems. I have just as little spoken of
Shakespeare's age. But when my reviewer in \textit{Dagbladet} at last asks what problems \textit{Lear} or
\textit{Macbeth} puts up for debate, I take the liberty, without further ado, of referring him, as a beginner in
aesthetic study, to the first handbook on Shakespeare that comes to hand, for example Ulrici's or Gervinus's. There
he will find an answer \dots{}\footnote{The work cited, p.~20.}

\medskip

\end{quote}

The war is carried into the enemy's own country. Rudolph Schmidt had reproached Georg Brandes for his imitations,
in content and language, of French literature:

\begin{quote}\small

What harm is there in that? When one walks in the sun, one gets sunburnt. But suppose now those accusations came
from a man of letters who has followed his teacher through thick and thin as a lackey follows his master, whose
opinion is the opinion Mr.~N. holds today, whose opinion was the one Mr.~N. held yesterday, and whose opinion will
be the one Mr.~N. will hold tomorrow! He forgets, then, to say that while I at least have breathed the free air of
Europe, he has breathed only his teacher's breath.\footnote{Ib. pp.~39 f. --- R. Schmidt wrote under the signature
``k''; Brandes, as one sees, breaks the pseudonymity and alludes to his faithful service as squire to ``N.'', i.e.\
Rasmus Nielsen.}

\medskip

\end{quote}

But at the end of each chapter he rises from the easily won bickering to general principles. This holds especially
of the polemic against ``44'', a pseudonym of Bishop Monrad. Finally he ends, like Brøchner in his reply to Rasmus
Nielsen, with a dignified emphasis on the fact that his conviction is firm and dearly bought, and therefore not
easily overturned by attacks.

\medskip

% --- p. 162 ---
\opage{162}\setcounter{footnote}{0}%
The above observations and reflections are, admittedly, only small matters. But when they concern an authorship
whose place in Danish literary history will quite certainly always be very prominent, they are not superfluous. It is
always worth the trouble to remember the names of the distinctive personalities who introduced the young Georg
Brandes to the problems of intellectual life, and to keep them in mind when reading his earliest writings.

\end{document}
